\newpage
\phantomsection
\label{prf:ordered-fields-are-densely-ordered}

\begin{remark*}[Return]
\hyperref[thm:ordered-fields-are-densely-ordered]{Return to Theorem}
\end{remark*}

\begin{theorem*}[Ordered Fields Are Densely Ordered]
Every ordered field is densely ordered by its order relation.
\end{theorem*}

\begin{proof}
Let $F$ be an ordered field. To show that $(F,<)$ is densely ordered, we verify
the two conditions in
\hyperref[def:densely-ordered-system]{Definition~\ref*{def:densely-ordered-system}}.

First, because $F$ is an ordered field, its order relation makes $F$ into an
ordered integral domain. In particular, for any $x,y\in F$, exactly one of the
relations
\[
x<y,\qquad x=y,\qquad y<x
\]
holds. Thus $(F,<)$ is a simply ordered system.

Now let $x,y\in F$ satisfy $x<y$. Define
\[
z:=\frac{x+y}{2}.
\]
Since $F$ is a field and $2=1+1\neq 0$, the element $2$ has a multiplicative
inverse, so $z\in F$.

Because $x<y$, the ordered-field axioms give
\[
x+x<x+y<y+y.
\]
Since $1>0$ by
\hyperref[thm:one-positive]{Theorem~\ref*{thm:one-positive}}, we have
$2=1+1>0$, and so
\hyperref[thm:inverse-positive]{Theorem~\ref*{thm:inverse-positive}} yields
$\frac12>0$. Multiplying the displayed inequality through by $\frac12$ and using
\hyperref[thm:multiply-positive]{Theorem~\ref*{thm:multiply-positive}}, we obtain
\[
\frac{x+x}{2}<\frac{x+y}{2}<\frac{y+y}{2}.
\]
By the field laws this simplifies to
\[
x<z<y.
\]
Therefore between any two distinct elements of $F$ there lies another element
of $F$. Hence $(F,<)$ is densely ordered.
\end{proof}

\begin{remark*}[Proof structure]
To prove density, one must exhibit a point strictly between two given ordered
field elements. The natural candidate is the midpoint $(x+y)/2$. The proof
first uses the ordered-field axioms to obtain $x+x<x+y<y+y$, then multiplies
through by the positive element $1/2$ to conclude
\[
x<\frac{x+y}{2}<y.
\]
This verifies the interpolation condition required for a densely ordered
system.
\end{remark*}

\begin{remark*}[Dependencies]~\\
\begin{itemize}
  \item \hyperref[def:densely-ordered-system]{Definition~\ref*{def:densely-ordered-system}}
  \item \hyperref[def:ordered-field]{Definition~\ref*{def:ordered-field}}
  \item \hyperref[thm:one-positive]{Theorem~\ref*{thm:one-positive}}
  \item \hyperref[thm:inverse-positive]{Theorem~\ref*{thm:inverse-positive}}
  \item \hyperref[thm:multiply-positive]{Theorem~\ref*{thm:multiply-positive}}
\end{itemize}
\end{remark*}
